\documentclass{beamer}
\mode<presentation> {
\usetheme{Warsaw}
\definecolor{myred}{RGB}{200,0,0} % Definir un nuevo color personalizado
\definecolor{moradoUPMH}{rgb}{0.46, 0.04, 0.49}

  % Cambiar el color de la estructura y de los elementos relacionados
  \setbeamercolor{structure}{fg=moradoUPMH}
  \setbeamercolor{palette primary}{fg=moradoUPMH, bg=white}
  \setbeamercolor{palette quaternary}{fg=white, bg=moradoUPMH}

  

  % Resto del código del tema

  %\colorlet{structure}{red} 
  %\usestructuretemplate{\color{structure}}{red}
  \beamertemplateshadingbackground{white!85}{moradoUPMH!15}
  \setbeamercovered{transparent}
}
\definecolor{darkblue}{rgb}{0.55,0,0}%{0,0,0.55}
\definecolor{springgreen4}{rgb}{0,0.55,0.27}
\definecolor{red4}{rgb}{0.55,0,0}
\definecolor{darkgreen}{rgb}{0,0.39,0}
\definecolor{darkred}{rgb}{0.55,0,0}
\definecolor{royalblue}{rgb}{0.55,0,0}%{0.23,0.37,0.8}
\definecolor{auburn}{rgb}{0.43, 0.21, 0.1}

\definecolor{verde}{rgb}{0,0.39,0}%{.5,.80,.5}
\definecolor{morado}{rgb}{.88,.50,.90}
\definecolor{azul}{rgb}{.5,.5,.98}
\definecolor{grisOscuro}{rgb}{.1,.1,.1}

\usepackage{subcaption}
\usepackage{tcolorbox}
\usepackage[table]{xcolor}
\usepackage{colortbl}
\usepackage{array, multirow, multicol, tabularx}

\usepackage{amsfonts}
\usepackage{amsmath}
\usepackage{amssymb}
\usepackage{graphics}
\usepackage{array}
\usepackage{xcolor}
\usepackage{setspace}
\usepackage{caption}
%\usepackage[spanish]{babel}
%%%%%%%%%%%%%\usepackage[spanish,es-tabla]{babel}
\usepackage{verbatim} % comentarios
\usepackage{multimedia}
\usepackage{media9}
\usepackage{ragged2e}
\usepackage{bm}




\usefonttheme[onlymath]{serif}
\usepackage{amsthm}
\usepackage{ae}
\usepackage{textcomp}
%\usepackage[utf8]{inputenc}
\usepackage{amssymb}
\usepackage{graphicx}
\usepackage[]{graphicx}
\newcommand{\bld}[1]{\mbox{$\mathbf{#1}$}}
%\newcommand{\e}{\bld{e}}
\newcommand{\cc}{\bld{c}}
\newcommand{\x}{\bld{x}}
\newcommand{\dx}{\dot{\bld{x}}}
\newcommand{\uu}{\bld{u}}
\newcommand{\nuu}{\bld{\nu}}
\newcommand{\vv}{\bld{v}}
\newcommand{\y}{\bld{y}}
\newcommand{\z}{\bld{z}}
\newcommand{\A}{\bld{A}}
\newcommand{\B}{\bld{B}}
\newcommand{\C}{\bld{C}}
\newcommand{\D}{\bld{D}}
\newcommand{\E}{\bld{E}}
\newcommand{\e}{\bld{e}}
\usepackage[outline]{contour}
%\spanishdecimal{.}
\setbeamertemplate{caption}[numbered] %%%%%%%Colocar para que ponga el numero de figura y tabla

\newcommand*\oldmacro{}%
\let\oldmacro\insertshorttitle%
\renewcommand*\insertshorttitle{%
\oldmacro\hfill%
\insertframenumber\,/\,\inserttotalframenumber}

\title[{August 25 - 27, 2025, San Diego, USA}]{Output Feedback Self-tuning PID Using Wavelet Neural Network Identification for Real-time UASs Trajectory Tracking}
\author[{9th IEEE CCTA 2025}]
{
    \texorpdfstring{
    %\vspace{-0.8cm}
     % \begin{columns}
     %            \column{6cm}
     %                \begin{flushleft}
     %                \includegraphics[scale = .07]{CongresoParis/Logo_IFAC_Paris.png}
     %                %\hspace{0.5cm}
     %                %\includegraphics[scale = .2]{LOGOS/cinve.png}
     %        \end{flushleft}
     %            \column{6cm}
     %              \begin{flushright}
     %                %\includegraphics[scale = .2]{LOGOS/UPMH-logo.png}
     %                %\includegraphics[scale = .5]{LOGOS/cinve.png}
     %         \end{flushright}
     %            \end{columns}
   %\vspace{-1.5cm}%Tenia -1.5
    %\\ Dr. Luis Rodolfo Garcia Carrillo \\
    % \vspace{0.5cm}
    % \begin{center}
    %     \includegraphics[scale = .07]{CongresoParis/Logo_IFAC_Paris_2.png}
    % \end{center}
    % MSc. Anel Olivares Cruz\\
    % Dr. Luis Enrique Ramos Velasco\\
    % Dr. Eduardo Steed Espinoza Quesada\\
    % MSc Omar Alejandro García Alcántara\\
    % Dr Luis Rodolfo Garcia Carrillo\\
    % Dr. Ignacio Rubio Scola
   Authors: \\
   \vspace{0.5cm}
       \begin{columns}
           \column{.5\linewidth}
           \centering
           MSc. Anel Olivares Cruz\\
           \\
           Dr. Eduardo S. Espinoza\\
           \column{.5\linewidth}
           \centering
           Dr. Luis E. Ramos Velasco\\
           \\
           \textbf{Dr. Luis R. Garcia Carrillo}\\
           \\
           %Dr. Ignacio Rubio Scola\\
			% \column{.23\linewidth}            
   %        \centering
   %         Omar Alejandro García Alcántara\\
       \end{columns}
   }
   { }
}
\date[July 2025]

\begin{document}
\contourlength{0.01cm} % Ajusta el grosor del contorno


	\begin{frame}
        \titlepage
	\end{frame}


    \begin{frame}
        \frametitle{Content}

        \tableofcontents[hideallsubsections]
        % \column{6cm}
        % \begin{block}{Next...}
        % \begin{center}
        %     \includegraphics [width=6cm,height=5cm]{Imagenes//Proyectar.jpg}
        % \end{center}
        % \end{block}

    \end{frame}

 
\section{Motivation}
\subsection{Motivation}
\begin{frame}{Motivation}
\begin{columns}
    
\column{5.5cm}
%Assumption:
\begin{table}[h]
    \centering
    \begin{tcolorbox}[
        tabularx = {l | l },
        title = Assumption:,
        colback = morado!10,   % Color de fondo 
        colframe = morado!80   % Color del borde
    ]
        $\star$ & Unknown mathematical \\
                 & model\\
        $\star$ & System identification and \\
        & control only with:\\
        &  input-output data\\
        % $\star$ & Investigación\\
    \end{tcolorbox}
\end{table}

\column{6cm}
\begin{block}{Unmanned Aerial Systems }
    \centering
    \includegraphics[width=\linewidth]{CongresoParis/frames_T(2).png}
\end{block}
    \end{columns}
\end{frame}
\subsection*{Experimental platform}


\begin{frame}{Experimental platform}

     \begin{figure}
         \centering
         \includegraphics[width=1.05\linewidth]{Imagenes San Diego/frames Neural-2.pdf}
     \caption{The schematic of the experimental flight testbed. The OptiTrack Motion Capture System (MCS) provides $\y[kT]$.}
         \label{fig:placeholder}
     \end{figure}

    %Fig. 1 
\end{frame}



  
\section{WaveNet-IIR Identification and Self-Tuning PID Controller}
\subsection{Diagram control block}
\begin{frame}{Control Scheme}

\begin{figure}[H]
    \centering
    \includegraphics[width=1.08\textwidth]{Imagenes San Diego/Control-Scheme-3.pdf}
    \raggedright
    \caption{Self-tuning PID WaveNet-IIR scheme}
    \label{fig:imagen_raggedright}
\end{figure}

\end{frame}



\subsection{(i) UAS: discrete input-output dynamics}

\begin{frame}{(i) UAS: discrete input-output dynamics}
%     State space representation of multiple-input multiple-output (MIMO) non-linear systems given by
% \begin{eqnarray}
% \label{xpc} \dx(t) & = & f(\x(t)) + g(\x(t))\uu(t)\\
% \label{ycon} \y(t) & =& h(\x(t)) \triangleq  \cc \x(t)
% \end{eqnarray}
% where $f(\x(t))$ is a vector field and $g(\x(t))$ is the input matrix.

The input and output of the discrete-time system  at the $(k+1)$ instant:
%
\begin{eqnarray}
\x[(k+1)T] = & \x[kT] + f(\x[kT])T + g(\x[kT])T\uu[kT]  \\
\y[(k+1)T]  = & \underbrace{\cc\left(\x[kT] + f(\x[kT])T\right)}_{\textcolor{blue}{\Phi[\x[kT],T]}}
+\underbrace{\cc(g(\x[kT])T)}_{\textcolor{blue}{\Gamma[\x[kT],T]}}\uu[kT]  \nonumber \label{out}
%&\triangleq &  \Phi[\x[kT],T] + \Gamma[\x[kT],T]\uu[kT]  \label{out}
\end{eqnarray}

 The strategy uses the control input $\uu[kT]$ and the UAS translational positions $\y[kT]$ to approximate the input-output mappings 
\begin{align}
    \hat \Gamma[kT] &= \Gamma[\x[kT],T]+e_\Gamma \nonumber\\ % and
    \hat \Phi[kT]       &= \Phi[\x[kT],T] +e_\Phi \\
    \hat{\y}[kT] &= \textcolor{blue}{\hat{\Phi}[kT] + \hat{\Gamma}[kT]}(u_1+u_2+u_3) \nonumber \label{estimacion}
\end{align}

%where $\epsilon_\Gamma$, and $\epsilon_\Phi$ are the finite approximation errors. From this procedure, the WaveNet-IIR filter yields
\end{frame}





\subsection{(ii) Identification (\textit{WaveNet-IIR})} 

% \begin{frame}{WaveNet-IIR identification}
%     The strategy uses the control input $\uu[kT]$ and the system output $\y[kT]$ (the UAS position translate) available to approximate the input-output mappings 
% \begin{align}
%     \hat \Gamma[\x[kT],T] &= \Gamma[\x[kT],T]+\epsilon_\Gamma \\ % and
%     \hat \Phi[\x[kT],T]       &= \Phi[\x[kT],T]] +\epsilon_\Phi
% \end{align}
% where $\epsilon_\Gamma$, and $\epsilon_\Phi$
% are the finite approximation errors. From this procedure, the WaveNet-IIR filter yields

% \end{frame}



\begin{frame}{(ii) Identification (\textit{WaveNet-IIR})}
    

    \begin{figure}
        \centering
            \includegraphics[width=0.9\columnwidth]{Imagenes San Diego/WaveNet-IIR scheme.pdf}\label{fig:Wavelet_neural_scheme}
        \caption{WaveNet-IIR filter has four layers: \textcolor{blue}{input}, \textcolor{darkgreen}{hidden}, \textcolor{red}{output}, and \textcolor{morado}{filtering} for the identification of the input-output UAS dynamics.} %Output $\textbf{z}$ is calculated by $\textbf{z}[kT]= \bm {f_n} [kT]W[kT]$.}
    \end{figure}

    % The synaptic weight matrices $W \in \mathbb{R}^{P\times J}$

\end{frame}

% \begin{frame}
% \justifying
%     The function $f_n (\tau)$ is called the "mother wavelet".The mother function is Morlet: \textcolor{blue}{$f_{n}(\tau)= \cos(\omega_0\tau)e^{-0.5\tau^{2}}$} \footnotemark
% \begin{align}
% \textcolor{blue}{f_{n,J} (\tau_J )=\frac{1}{\sqrt{(s_J )}} f_n (\tau_J ), \tau_J (k,q,s)=\frac{k-q_J}{s_J}}
% \end{align}
% Where $q$ and $s$ shift and scale the signal, with respect to $k$ indicating the moment of the operation and $\tau$ being a discrete variable in time.\\

% \begin{tcolorbox}[
%         tabularx = {l | c},
%         %title = Tabla,
%         colback = morado!10,   % Color de fondo 
%         colframe = morado!80   % Color del borde
%     ]
%         Parameters  & Matrix form \\
%         \hline \hline
%        Translation  &$\bm{q}=[q_1,q_2,\ldots,q_J]^T$\\
%        Dilation & $\bm {s}=[s_1,s_2,\ldots,s_J ]^T$\\
%        "Daughter" wavelets & $\bm {f_n} =[f_{n,1},f_{n,2},\ldots,f_{n,J} ]^T$
%     \end{tcolorbox}

% \vspace{0.1cm}%Tenia -1.5

%     \tiny{\footnotemark[1]Domínguez Mayorga, C., Espejel Rivera, M., Ramos Velasco, L., Ramos Fernández, J. y Escamilla Hernández, E. (2012) "Wavenet Algorithms with Applications in Approxima-tion Signals: A Comparative Study.", Revista Iberoamericana de Automática e Informática industrial, 9(4), pp. 347–358. doi: 10.1016/j.riai.2012.09.001.
%     }
% \end{frame}




\begin{frame}{(ii) Identification (\textit{WaveNet-IIR}): IIR Filter}

\begin{figure}
        \centering
            \includegraphics[width=1.1 \linewidth]{Imagenes San Diego/Escructura del filtroIIR.pdf}
        \caption{IIR filter with $\bm{A}$ and $\bm{B}$ representing the number of feedforward and feedback coefficients, respectively.}
    \end{figure}
\end{frame}



% \begin{frame}{WaveNet-IIR Neural-network}
% \justifying
% With the help of IIR filters, the estimated positions $\hat{\textbf{y}}$ can be calculated, that is, the system output-input is described by:
% \begin{align}
% \textcolor{blue}{\hat{\textbf{y}}[k]} &= \textcolor{blue}{\hat{\Phi}[k] + \hat{\Gamma}[k] U}
% \label{estimacion}
% \end{align}

%     \begin{tcolorbox}[
%         tabularx = {l  },
%         %title = Tabla,
%         colback = morado!10,   % Color de fondo 
%         colframe = morado!80   % Color del borde
%     ]
%         here:  \\
%         \hline \hline
%         $
% \hat{\Phi}[k] =
% \begin{bmatrix}
% B \cdot Y^\top
% \end{bmatrix} s_p[k]$\hspace{4.5cm}$B\in \mathbb{R}^{P\times R}$\\

% $\hat{\Gamma}[k] =
% \begin{bmatrix}
% A \cdot Z^\top
% \end{bmatrix}$ \hspace{5.2cm} $A \in \mathbb{R}^{P \times Q}$\\

% $
% Z =
% \begin{bmatrix}
% \textbf{z}[k],
% \textbf{z}[k - 1],
% \dotsb,
% \textbf{z}[k - Q]
% \end{bmatrix}$\\

% $Y[k] =
% \begin{bmatrix}
% \hat{\textbf{y}}[k - 1],
% \hat{\textbf{y}}[k - 2],
% \dotsb,
% \hat{\textbf{y}}[k - R]
% \end{bmatrix}
% $\\ 

% $U = u_1[k] + u_2[k] + u_3[k] $
%     \end{tcolorbox}


% \end{frame}



% \begin{frame}{Estimation cost function}
%     %The estimation cost function is given by
% \begin{align}
% E[kT] = \frac{1}{2} \sum_{i} e_i^2[kT]
% \end{align}
% with the estimation error given by $ e_i[kT] = \textbf{y}_i[kT] - \hat{\textbf{y}}_i[kT]$. Here, \( \textbf{y}_i \) represents each output of the system and \( \hat{\textbf{y}}_i \) is the estimated output of the WaveNet-IIR, for $ i = (x, y, z)$.
% \end{frame}




\begin{frame}{(ii) Identification (\textit{WaveNet-IIR}): Gradient descent}

The estimation cost function $E[kT] = \frac{1}{2} \sum_{i} e_i^2[kT]$ can be minimized making use of the gradient descent method: %, which uses the set of partial derivatives given by
\begin{align}
\frac{\partial E}{\partial \bm {W}[kT]} &=  \textcolor{blue}{UI_e(\bm {A}f_a)^\top} \\
\frac{\partial E}{\partial \bm {q}[kT]} &=  \textcolor{blue}{Ue(\bm {A}f_b)^\top} \\
\frac{\partial E}{\partial \bm {s}[kT]} &=  \textcolor{blue}{\frac{\partial E}{\partial \bm {q}[kT]} \odot  \tau} \\
\frac{\partial E}{\partial \bm {A}[kT]} &=  \textcolor{blue}{UI_eZ^\top} \\
\frac{\partial E}{\partial \bm {B}[kT]} &= \textcolor{blue}{s_pI_e(\hat{Y}f_a)^\top}
\end{align}
\end{frame}

% \begin{frame}{Gradient descent method}
%     where
% \begin{align}
% I_e     &= \text{diag}( e_x[kT], e_y[kT], e_z[kT]) \\
% e[kT]   &= [e_x[kT] e_y[kT] e_z[kT] ] \\
% f_a[kT] &=[ \bm {f_n}[kT], \bm {f_n}[k-1], \ldots, \bm {f_n}[k-Q] ]^\top \\
% f_b[kT] &= \frac{\partial f_a[kT]}{\partial \tau}
% \end{align}
% \end{frame}

\section{Design of the self-tuning PID control strategy.}
%\subsection{Classic discrete PID control law}
% \begin{frame}{Discrete PID control law}

% \begin{align}
%     \label{eq:control law}
%     \textcolor{verde}{\mathbf{u}[k+1]} &= \textcolor{verde}{\mathbf{u}[k] + \bm{K}_p[k][\boldsymbol{\epsilon}[k] - \boldsymbol{\epsilon}[k-1]] + \bm{K}_i[k][\boldsymbol{\epsilon}[k]]} \nonumber \\ 
%     & \textcolor{verde}{+ \bm{K}_d[k][\boldsymbol{\epsilon}[k] - 2\boldsymbol{\epsilon}[k-1] + \boldsymbol{\epsilon}[k-2]] }
% \end{align}
   

% \begin{table}[h]
%     \centering
%     \begin{tcolorbox}[
%         tabularx = {l | l },
%         %title = Tabla,
%         colback = morado!10,   % Color de fondo 
%         colframe = morado!80   % Color del borde
%     ]
%         where: & \\
%         \hline \hline
%         $\bm{K}_p$ & proportional gain, $\bm{K}_p= \text{diag}\{K_{p_x},K_{p_y},K_{p_z}\}$\\
%         $\bm{K}_i$ & integral gain, $\bm{K}_i=\text{diag}\{K_{i_x},K_{i_y},K_{i_z}\}$\\
%         $\bm{K}_d$ & derivative gain, $\bm{K}_d=\text{diag}\{K_{d_x},K_{d_y},K_{d_z}\}$\\
%         $\epsilon$ & tracking error $\epsilon=\boldsymbol{\xi}-\boldsymbol{\xi}_{ref}$\\
%         $\boldsymbol{\xi}$ & represents the UAS's position $\boldsymbol{\xi}=(x, y,z)^\top$\\
%         $\boldsymbol{\xi}_{ref}$ & desired references $\boldsymbol{\xi}_{\text{ref}}=(x_{\text{ref}}, y_{\text{ref}},z_{\text{ref}})^\top$\\
%     \end{tcolorbox}

% \end{table}

%\end{frame}

\subsection{(iii) PID controller and (iv) Auto-tuning gains}
 
\begin{frame}{(iii) PID controller and (iv) Auto-tuning gains}
\vspace{-0.5cm}
\begin{align}
    \label{eq:control law}
    \textcolor{verde}{\mathbf{u}[(k+1)T]} &= \textcolor{verde}{\mathbf{u}[kT] + \bm{K}_p[kT][\boldsymbol{\epsilon}[kT] - \boldsymbol{\epsilon}[(k-1)T]] + \bm{K}_i[kT][\boldsymbol{\epsilon}[kT]]} \nonumber \\ 
    & \textcolor{verde}{+ \bm{K}_d[kT][\boldsymbol{\epsilon}[kT] - 2\boldsymbol{\epsilon}[(k-1)T] + \boldsymbol{\epsilon}[(k-2)T]] }
\end{align}
   These gains can be dynamically tuned following the cost function for the tracking error $\mathcal{E}[kT] = \frac{1}{2} \sum_i \epsilon_i^2[kT]$.
% \begin{equation}
% \mathcal{E}[kT] = \frac{1}{2} \sum_i \epsilon_i^2[kT]
% \end{equation}
Applying the partial derivative for each gain:
\begin{align*}
        \nabla \bm{K}_p[kT] &=  \textcolor{auburn}{e[kT] \hat{\Gamma}[kT](\epsilon[kT]-\epsilon[(k-1)T])}\\
        \nabla \bm{K}_i [kT] &=  \textcolor{auburn}{e[kT] \hat{\Gamma}[kT](\epsilon[kT])}\\
        \nabla \bm{K}_d [kT] &=  \textcolor{auburn}{e[kT] \hat{\Gamma}[kT](\epsilon[kT]-2 \epsilon[(k-1)T]  + \epsilon[(k-2)T])}
\end{align*} 
The parameter update rule:
\begin{align}
\bm{K}_{*} [(k+1)T] &= \bm{K}_{*} [kT] + \mu_{K_{*}} \nabla \bm{K}_{*} [kT]
\end{align}

\end{frame}

% \begin{frame}{Gradient descent}
% Applying the partial derivative for each gain we obtain
% \begin{align}
%         \nabla \bm{K}_p &=  \textcolor{auburn}{e[kT] \hat{\Gamma}[kT](\epsilon[kT]-\epsilon[k-1])}\label{autotunning3}\\
%         \nabla \bm{K}_i &=  \textcolor{auburn}{e[kT] \hat{\Gamma}[kT](\epsilon[kT])} \label{autotunning2} \\
%         \nabla \bm{K}_d &=  \textcolor{auburn}{e[kT] \hat{\Gamma}[kT](\epsilon[kT]-2 \epsilon[k-1]  + \epsilon[k-2])}
%         \label{autotunning}
% \end{align} 
% where the input matrix \textcolor{blue}{$\hat{\Gamma}$} is obtained from the \textit{WaveNet}-IIR.%the identification process performed by

% Finally, the parameter update rule is expressed as
% \begin{align}
% K_{*} [k+1] &= \textcolor{auburn}{K_{*} [kT] + \mu_{K_{*}} \frac{\partial \mathcal{E}[KT]}{\partial K_{*}[KT]}}
% \end{align}
% where $\mu_{K_{*}}$ is the learning rate.
% \end{frame}





\section{Experimental results}

\subsection{Training of the WaveNet-IIR}


\begin{frame}{Parameter configuration of the \textit{WaveNet}-IIR}
%The parameter configuration of the \textit{WaveNet}-IIR, designed to identify the translational dynamics of the quad rotorcraft using the mother wavelet is presented in Table~\ref{tab:init config morlet}.
\small{
\begin{table}[h!]
        \centering
        \caption{Configuration of the \textit{WaveNet}-IIR.}
        \label{tab:init config morlet}
        \begin{tabular}{cc}
            \hline
            \textbf{Parameter} & \textbf{Value}  \\ \hline   \hline 
           Number of neurons, $J$ & 5 \\
             Number of inputs, $I$ & 3\\
             Number of outputs, $P$  & 3\\\
             Lead coefficients, $R$ & 3\\
             Lag coefficient, $Q$  & 4 \\\
             Sampling frequency, $1/T$  & 100 Hz\\
             Sampling period, $T$&  $10ms$\\
             Persistent signal, $s_p$     & $[5, 17, 17]$\\
             \hline   \hline
        \end{tabular}
    \end{table}}
   
\end{frame}

\begin{frame}{Initial values}
    % From trained tests, the initial operational parameters presented in Table~\ref{tab:init values morlet} were determined.
\small{
\begin{table}[H]
        \centering
        \caption{Initial values for the \textit{WaveNet}-IIR.}
        \label{tab:init values morlet}
        \begin{tabular}{cc}
            \hline
            \textbf{Parameter} & \textbf{Value} \\ \hline   \hline
            $\bm{W}[kT]$ &  $
                    \begin{bmatrix}
                    -0.009 & -0.014 & -0.004 & -0.009 & -0.013 \\
                    0.016 & 0.023 & 0.024 & 0.028 & 0.017 \\
                    -0.050 & -0.038 & -0.036 & -0.041 & -0.041
                    \end{bmatrix}
                    $ \\ \vspace{.1 cm}
            $\bm{A}[kT]$ &   $
                    \begin{bmatrix}
                    -0.030 & -0.030 & 0.012 & 0.030 & 0.013 \\
                    0.006 & 0.005 & 0.005 & 0.006 & 0.007 \\
                    -0.018 & -0.012 & 0.003 & 0.006 & 0.001
                    \end{bmatrix}
                    $ \\ \vspace{.1 cm}
            $\bm{B}[kT]$ &  $
                    \begin{bmatrix}
                    0.025 & 0.083 & 0.092 \\
                    0.006 & 0.022 & 0.030 \\
                    0.011 & 0.026 & 0.028
                    \end{bmatrix}
                    $ \\ \vspace{.1 cm}
            $\bm{s}[kT]$ &   $\begin{bmatrix} 4, & 4, & 4, &  4, & 4\end{bmatrix} $ \\ \vspace{.1 cm}
            $\bm{q}[kT]$ &   $\begin{bmatrix}  1, &  2, &  3, & 4, & 5\end{bmatrix} $ \\
            \hline   \hline
        \end{tabular}
    \end{table}
}
\end{frame}


% \subsection{Rhodonea Curve Trajectory}
% \begin{frame}{Rhodonea Curve Trajectory}
%     The UAS starts tracking the Rhodonea curve, which resembles a four-petal rose shape on the \(x_{I}-y_{I}\) plane. %The parameters describing this particular rose shape are given by:
%     \begin{align}
%         x_{\text{ref}} (nT_{\text{neu}} ) = r_c \cos{\left(l  \omega n T_{\text{neu}}  \right)} \cos{\left(\omega n T_{\text{neu}} \right)}\\
%          y_{\text{ref}}\left( n T_{\text{neu}}  \right) = r_c \cos{\left(l \omega n T_{\text{neu}} \right)} \sin{\left(\omega n T_{\text{neu}} \right)}
%     \end{align}
    
%     \begin{table}[h]
%     \centering
%     \begin{tcolorbox}[
%         tabularx = {l | l },
%         %title = Tabla,
%         colback = morado!10,   % Color de fondo 
%         colframe = morado!80   % Color del borde
%     ]
%         where: &\\
%         \hline \hline
%         $n$ & is the n-th sample\\
%         \(r_c = 1\) & specify the length of petals\\
%         \(l = 2\) & specify the number of petals\\
%        \(\omega = \pi/25\) & is the angular frequency\\

%     \end{tcolorbox}
%     \end{table}

% \end{frame}








\begin{frame}{Initial values}

\begin{table}[H]
        \centering
        \caption{Initial values: learning rate and self-tuning PID.}
        \label{tab: init Param PIDs}
        \begin{tabular}{cccc}
            \hline
            \textbf{Learning rate} & \textbf{Value} & \textbf{Ganis} & \textbf{Value}  
            
            \\ \hline   \hline 
            $\mu_\mathbf{W}$ & $2.5\times10^{-2}$&  $\bm{K}_i[kT]$ & $\text{diag}\{ 0.2, \, 0.2, \, 1 \} $ \\
            $\mu_\mathbf{q}$ & $1$&  $\bm{K}_d[kT]$ & $\text{diag}\{ 9, \, 9, \, 15 \} $ \\
             $\mu_\mathbf{s}$ & $1$&$\bm{K}_p [kT]$ & $\text{diag}\{ 50, \, 50, \, 100 \} $\\
             $\mu_\mathbf{B}$ & $1.6\times10^{-3}$&  & \\
             $\mu_\mathbf{A}$ & $7\times10^{-3}$&  & \\
             $\bm{\mu}_{\bm{K}_{d}} $ & $ [ 8, \,8,\, 3] $&  & \\
            $\bm{\mu}_{\bm{K}_{i}}$ & $[ 0.3, \, 0.3,\, 0.1 ] $& & \\  
            $\bm{\mu}_{\bm{K}_{p}}$&$ [ 30, \, 30, \, 30 ] $& &\\ 
            %   & &
            \hline \hline
        \end{tabular}
    \end{table}
\end{frame}

\subsection{Self-tuning PID control: trajectory tracking results}

\begin{frame}{Experimental results of trajectory tracking}
\begin{figure}
        \centering
            \includegraphics[width=0.72 \linewidth]{Imagenes San Diego/MovimientoUAV-1.pdf}
        \caption{The real-time tracking of the\href{https://youtu.be/vWSzDwxYs_8}{ Rhodonea curve} in a 3D space .}
    \end{figure}
     

\end{frame}

\begin{frame}{Experimental results of trajectory tracking}

\begin{figure}
        \centering
             \includegraphics[width=0.85 \linewidth]{Imagenes San Diego/Graficas_xyz_sinEst-1.pdf}
        \caption{Tracking performance using PID \textit{WaveNet}-IIR.}
    \end{figure}

\end{frame}


\subsection{Self-tuning PID control: model identification results}

\begin{frame}{Experimental results of trajectory identification }

\begin{columns}
    \column{7cm}
    \begin{figure}
        \centering
             \includegraphics[width=0.9 \linewidth]{Imagenes San Diego/GraficaXYZEstimacion.pdf}
        \caption{Output identification using the \textit{WaveNet}-IIR.}
    \end{figure}

    \column{6 cm}
    \vspace{-1cm}
    \small{
    \begin{table}[h!]
        \centering
        \caption{Evaluation parameters of the \textit{WaveNet}-IIR}
        \label{tab:estimation error}
        \begin{tabular}{ccc}
            \hline
            \textbf{Output} & \textbf{RMSE} & \textbf{Correlation}\\ \hline   \hline
            $z$ &  0.057 & 0.995 \\ %\vspace{.1 cm}
            $x$ &   0.110 & 0.971\\ %\vspace{.1 cm}
            $y$ &  0.165 & 0.901\\
            \hline   \hline
        \end{tabular}
    \end{table}}
\end{columns}

\end{frame}

% \begin{frame}{Frame Title}
%     The small RMSE values and correlations exceeding 0.9, detailed in Table~\ref{tab:estimation error}, demonstrate the \textit{WaveNet}-IIR's high estimation accuracy.

% \begin{table}[h!]
%         \centering
%         \caption{Root Mean Square Error (RMSE) and correlation of the \textit{WaveNet}-IIR}
%         \label{tab:estimation error}
%         \begin{tabular}{ccc}
%             \hline
%             \textbf{Output} & \textbf{RMSE} & \textbf{Correlation}\\ \hline   \hline
%             $z$ &  0.057 & 0.995 \\ %\vspace{.1 cm}
%             $x$ &   0.110 & 0.971\\ %\vspace{.1 cm}
%             $y$ &  0.165 & 0.901\\
%             \hline   \hline
%         \end{tabular}
%     \end{table}
% \end{frame}




\section{Conclusions}

\begin{frame}{Conclusions}
    
    

    % \begin{columns}
    %             \column{8.5cm}
                \begin{itemize}
                \item The WaveNet-IIR proved effectiveness in the identification of the UAS dynamics.
                \item The WaveNet-IIR capability for real-time trajectory tracking, with self-tuning gains, was successfully demonstrated.
                 \end{itemize}
    %             \column{3.5cm}
    %             \begin{block}
    %             \begin{center}
    %             \includegraphics [width=3.5cm,height=5cm]{Imagenes//goldfigure.jpg}
    %             \end{center}
    %             \end{block}
    %         \end{columns}
\end{frame}


        \begin{frame}
            %\frametitle{Aclaraciones, Dudas y Comentarios}
             
       
                \begin{center}
                \textbf{\Huge Thank you for your attention}\\
                luisillo@nmsu.edu \\
                \end{center}

        \end{frame}
	
	
	
\end{document}